\paragraph{Instancias.} El experimento se realizó sobre 288 instancias generadas con \texttt{grid\_generator.py} (semilla 2026), abarcando seis valores de filas $N \in \{4, 8, 16, 64, 256, 1024\}$ y dos proporciones de columnas $M \in \{N, 2N\}$ (12 tamaños de grilla desde $4\times 4$ hasta $1024\times 2048$). Para cada tamaño se generaron cuatro tipos de terreno (\texttt{libre}, \texttt{murallas}, \texttt{minas} y \texttt{mixta}), dos dominios de valores (\texttt{D10} con penalización $P=10$ y \texttt{D1000} con $P=1000$) y tres muestras por combinación ($a,b,c$), fijando $V=3$ vidas iniciales. El generador reserva un camino monótono libre de obstáculos para garantizar que toda instancia tenga al menos una solución válida. Fuerza bruta se ejecutó únicamente cuando $N+M \le 32$ (120 de las 288 instancias), generando un total de 696 archivos de salida entre los tres algoritmos.

\paragraph{Implementación.} Los tres algoritmos se implementaron en C++17 (\texttt{g++ -O2}) retornando una estructura \texttt{PathResult} con el puntaje obtenido, la secuencia de movimientos (\texttt{R}/\texttt{D}) y las vidas restantes, o bien indicando que no se encontró solución. Fuerza bruta explora los caminos mediante \emph{backtracking} descartando solo ramas inválidas (murallas o vidas agotadas); programación dinámica resuelve el problema de forma \emph{bottom-up} con estado $(i,j,k)$ según las $k$ minas detonadas y reconstruye el camino óptimo; y greedy aplica en cada paso la regla local del enunciado, avanzando hacia la celda vecina accesible de mayor ganancia neta (desempatando hacia la derecha).

\paragraph{Mediciones.} Los tiempos se registraron con \texttt{std::chrono::steady\_clock} (1 ejecución para fuerza bruta y el mínimo de 5 repeticiones para programación dinámica y greedy). El consumo de memoria corresponde al \emph{pico de heap adicional} medido mediante sobrecarga de \texttt{operator new/delete}, excluyendo la grilla de entrada. Todas las ejecuciones corrieron de manera secuencial en la misma máquina.

\paragraph{Verificación.} El programa principal re-simula cada camino devuelto para validar que no atraviesa murallas, que no agota las vidas y que reproduce exactamente el puntaje y las vidas reportadas; además comprueba que fuerza bruta y programación dinámica obtengan el mismo puntaje en las 120 instancias donde ambos corren, y que $s(C_{\text{greedy}})\le s(C^*)$. Todas las verificaciones finalizaron con 0 errores, tomando el puntaje de programación dinámica como referencia óptima para $\Delta = s(C^*) - s(C_{\text{greedy}})$.